%# -*- coding: utf-8-unix -*-

\begin{thebibliography}{99}
	\setlength{\itemsep}{3pt plus 1pt minus 1pt}   % 黄金比例紧凑纵向间距
	\setlength{\parsep}{1pt plus 0.5pt}            % 紧凑的段落空隙
	\small                                         % 学术出版级别小五号字

	\bibitem{Borevich} BOREVICH, Z. I. and SHAFAREVICH, I. R., \textit{Number theory} (Moscow, 1964; Academic Press, London, 1966).

	\bibitem{Cassels1} Cassels, J. W. S., \textit{An introduction to Diophantine approximation} (Cambridge Univ. Press, 1957).

	\bibitem{Cassels2} Cassels, J. W. S., \textit{An introduction to the geometry of numbers} (G{\"o}ttingen, Berlin; Springer Verlag, Heidelberg, First edn. 1959).

	\bibitem{Feldman} FELDMAN, N. I. and Shidlovsky, A. B., The development and present state of the theory of transcendental numbers, \textit{Russian Math. Surveys} 22 (1967), 1--79; translated from \textit{Uspehi Mat. Nauk SSSR}, 22 (1967), 3--81.

	\bibitem{Gelfond1} GELFOND, A. O., \textit{Transcendental and algebraic numbers} (Moscow, 1952; Dover publ., New York, 1960).

	\bibitem{Gelfond2} GELFOND, A. O., \textit{Selected works} (Moscow, 1973).

	\bibitem{Gelfond3} GELFOND, A. O. and LINNIK, Yu. V., \textit{Elementary methods in the analytic theory of numbers} (Moscow, 1962; Pergamon Press, Oxford, 1966).

	\bibitem{Hecke} HECKE, E., \textit{Theorie der algebraischen Zahlen} (Teubner, Leipzig, 1923).

	\bibitem{Koksma} Koksma, J. F., \textit{Diophantische Approximationen} (Ergebnisse Math., Berlin and Leipzig, 1936).

	\bibitem{Lang1} LANG, S., \textit{Introduction to transcendental numbers} (Addison-Wesley, Reading Mass., 1966).

	\bibitem{Lang2} LANG, S., \textit{Diophantine geometry} (Interscience, New York, 1962).

	\bibitem{LeVeque} LEVEQUE, W. J., \textit{Topics in number theory} (Addison-Wesley, Reading Mass., 1956).

	\bibitem{Mahler} MAHLER, K., \textit{Lectures on Diophantine approximations} (Notre Dame, 1961).

	\bibitem{Maillet} MAILLET, E., \textit{Introduction {\`a} la th{\'e}orie des nombres transcendants et les propri{\'e}t{\'e}s arithm{\'e}tiques des fonctions} (Gauthier-Villars, Paris, 1906).

	\bibitem{Niven} NIVEN, I., \textit{Numbers, rational and irrational} (Interscience, New York, 1961).

	\bibitem{Perron} Perron, O., \textit{Die Lehre von den Kettenbr{\"u}chen} (Teubner, Leipzig, 1913).

	\bibitem{Ramachandra} RAMACHANDRA, K., \textit{Lectures on transcendental numbers} (Ramanujan Inst., Univ. Madras, 1969).

	\bibitem{Schmidt} SCHMIDT, W. M., \textit{Lectures on Diophantine approximation} (Dept. Math., Univ. Colorado, 1970).

	\bibitem{Schneider} Schneider, Th., \textit{Einf{\"u}hrung in die transzendenten Zahlen} (G{\"o}ttingen, Berlin; Springer Verlag, Heidelberg, 1957).

	\bibitem{Siegel} SIEGEL, C. L., \textit{Transcendental numbers} (Princeton, 1949).

	\bibitem{Skolem} Skolem, Th., \textit{Diophantische Gleichungen} (Ergebnisse Math., Berlin and Leipzig, 1938; reprinted: Chelsea publ., New York, 1950).

	\bibitem{Sprindzuk1} SPRIND{\v{Z}}UK, V. G., \textit{Mahler's problem in metric number theory} (Minsk, 1967; American Math. Soc. Transl. Math. Monographs, vol. 25, 1969).

	\bibitem{Sprindzuk2} SPRIND{\v{Z}}UK, V. G. (Editor), \textit{Current problems in the analytic theory of numbers} (Minsk, 1974).

	\bibitem{Waldschmidt} WALDSCHMIDT, M., \textit{Nombres transcendants}, Lecture Notes in Mathematics, vol. 402 (Springer-Verlag, Berlin, 1974).

\end{thebibliography}
